\section{Related Work}
\label{sec:related}

\paragraph{VLA policies and their generalization.}
VLAs fine-tune vision-language backbones to output robot actions, as discrete tokens~\cite{brohan2023rt2,kim2024openvla}
or through a flow-matching action expert attached to the backbone~\cite{black2024pi0,pi2025pi05}; continuous-action
fine-tuning brought them close to saturation on LIBERO~\cite{liu2023libero,kim2025oft}.
Robustness benchmarks built on LIBERO show how narrow this success is: LIBERO-Plus~\cite{fei2025liberoplus} perturbs
camera, lighting, robot initial state and layout, and LIBERO-PRO~\cite{zhou2025liberopro} swaps and moves objects,
rephrases instructions and changes tasks, with near-zero success for most models under position changes.
We study the position failures of \pizf, the strongest open model on these suites.

\paragraph{Memorization and shortcuts in imitation.}
Behavior cloning is known to latch onto spurious correlates of the expert's action, such as past actions or
proprioception~\cite{dehaan2019causal,wen2020copycat}, an instance of shortcut learning~\cite{geirhos2020shortcut}.
For VLAs, recent analyses report that instructions retrieve memorized trajectories~\cite{chen2026ect}, that new target
positions remain decodable from internal features while the arm goes to the old one~\cite{encoded2026}, and that spatial
motor programs are tied to the scene layout~\cite{features2026}.
VLA-Trace~\cite{shi2026vlatrace} traces representations and attention pathways of \pizf and OpenVLA during fine-tuning.
A related line studies \emph{instruction} blindness, where visual priors override language: LIBERO-CF and counterfactual
action guidance~\cite{fang2026whenvision}, deconfounding with a language-masked branch~\cite{zhang2026cofactvla}, residual
semantic steering~\cite{zhan2026stable}, and flatness-preserving fine-tuning~\cite{flatness2026}.
Our measurements agree with the memorization diagnosis; what we add is a per-phase, per-intervention separation of use and
encoding that compares models before and after a remedy, and a localization of the parameters whose change restores use.

\paragraph{Adding geometry.}
Many works feed depth, point clouds or 3D features to VLAs~\cite{zhen20243dvla,qu2025spatialvla}; where position
perturbations were measured, they did not resolve them.
Closest to our question, Tsai \etal~\cite{tsai2026point} inject a grounded 3D goal point into the adaptive normalization of
the action expert and report large gains when training on 69 tasks with varied layouts.
We re-implement the injection with an oracle point in our per-suite setting and find that the channel is ignored; with
counterfactual supervision added, the policy still solves the task through the VLM rather than through the point.

\paragraph{Escaping memorized trajectories.}
Inference-time methods leave the policy unchanged: Affordance Field Intervention~\cite{xu2026afi} detects ``memory traps''
and steers \pizf with 3D affordance fields and waypoints (reporting +20.2\% on LIBERO-PRO), and CoRe~\cite{core2026} plans
recoveries in imagination. CounterAlign~\cite{counteralign2026} learns a reward from counterfactual discriminators for
offline RL and reports gains under position perturbations. These are complementary to ours: they act at test time or
through a learned reward, whereas we ask where in the policy the failure lives and supervise that decision directly.

\paragraph{Counterfactual data and privileged teachers.}
Counterfactual data augmentation~\cite{pitis2020coda} and demonstration generation by replaying transformed
trajectories~\cite{mandlekar2023mimicgen,xue2025demogen} create new supervision from existing episodes.
\ect~\cite{chen2026ect} is the closest method: it mirrors or shifts the whole scene at the start of a demonstration and
replays the transformed end-effector path, training on both.
Our counterfactuals are instead built at \emph{arbitrary states of on-policy rollouts}, one change at a time (the target
before the grasp, the container while carrying), and labeled by a privileged scripted teacher, in the spirit of learning
from privileged experts~\cite{chen2019lbc} and on-policy distillation~\cite{ross2011dagger,agarwal2024onpolicy}.
We use them as an instrument to ask where VLAs fail, and treat the resulting training procedure as the minimal
intervention the analysis calls for rather than as a data-generation contribution.
